\begin{abstract}
The expressivity of artificial neural networks is classically studied through approximation theory, which provides asymptotic guarantees but obscures the exact algebraic boundaries of finite architectures. Motivated by the adoption of polynomial activation functions in niche aspects of machine learning, we reframe neural network expressivity as a problem in algebraic geometry. By restricting activations to quadratic forms, the network's parameterized function space maps to a secant variety of a Veronese embedding. We apply elimination theory to explicitly derive the implicit vanishing ideal for an example architecture. Our results demonstrate both the existence of universally expressive architectures and also precise algebraic limitations for others. This algebraic framework provides an exact certificate of learning capabilities, independent of training data or optimization heuristics.
\end{abstract}